\documentclass{article}
\usepackage{amsmath,amssymb,amsthm}
\usepackage{algorithm}
\usepackage{algpseudocode}
\usepackage{enumitem}
\newtheorem{theorem}{Theorem}
\newtheorem{lemma}{Lemma}
\newtheorem{assumption}{Assumption}
\newtheorem{remark}{Remark}
\begin{document}

\section*{Proximal Gradient Method (PGM)}

\section*{Mathematical Formulation}
Let  
\[
F(x):=f(x)+g(x),\qquad f:\mathbb{R}^{d}\to\mathbb{R}\ \text{convex, $L$-smooth},\qquad   
g:\mathbb{R}^{d}\to\mathbb{R}\cup\{+\infty\}\ \text{convex, possibly non‑smooth}.
\]
For a stepsize $\eta>0$, the proximal gradient iteration is  
\[
\boxed{\,x_{k+1}= \operatorname{prox}_{\eta g}\bigl(x_{k}-\eta\nabla f(x_{k})\bigr)\,},\qquad k=0,1,\dots
\]
where the proximal operator of $g$ is  
\[
\operatorname{prox}_{\eta g}(v):=\arg\min_{u\in\mathbb{R}^{d}}\Bigl\{ g(u)+\frac{1}{2\eta}\|u-v\|^{2}\Bigr\}.
\]

\section*{Pseudocode}
\begin{algorithm}[H]
\caption{Proximal Gradient Method}
\begin{algorithmic}[1]
\State \textbf{Input:} stepsize $\eta\in(0,1/L]$, initial point $x_{0}\in\mathbb{R}^{d}$, max iter $K$.
\For{$k=0$ \textbf{to} $K-1$}
    \State $v_{k}\gets x_{k}-\eta\nabla f(x_{k})$ \Comment{gradient step}
    \State $x_{k+1}\gets \operatorname{prox}_{\eta g}(v_{k})$ \Comment{proximal step}
\EndFor
\State \textbf{Output:} $x_{K}$
\end{algorithmic}
\end{algorithm}

\section*{Dry Run Example}
Consider the scalar problem  
\[
f(w)=(w-3)^{2},\qquad g\equiv 0,\qquad \eta=0.1,\qquad w_{0}=0 .
\]
Since $g\equiv0$, $\operatorname{prox}_{\eta g}(v)=v$.

\begin{enumerate}[label=\textbf{Iter \arabic*:},leftmargin=*]
\item $k=0$  
\[
\nabla f(w_{0})=2(w_{0}-3)=-6,\qquad 
v_{0}=w_{0}-\eta\nabla f(w_{0})=0-0.1(-6)=0.6,
\]
\[
w_{1}=v_{0}=0.6,\qquad F(w_{1})=(0.6-3)^{2}=5.76.
\]

\item $k=1$  
\[
\nabla f(w_{1})=2(0.6-3)=-4.8,\qquad 
v_{1}=0.6-0.1(-4.8)=1.08,
\]
\[
w_{2}=v_{1}=1.08,\qquad F(w_{2})=(1.08-3)^{2}=3.6864.
\]

\item $k=2$  
\[
\nabla f(w_{2})=2(1.08-3)=-3.84,\qquad 
v_{2}=1.08-0.1(-3.84)=1.464,
\]
\[
w_{3}=v_{2}=1.464,\qquad F(w_{3})=(1.464-3)^{2}=2.365.
\]
\end{enumerate}
The objective value decreases monotonically, illustrating the algorithm’s progress.

\newpage

\section*{Assumptions}
\begin{assumption}[Smoothness of $f$]\label{ass:smooth}
$f$ is convex and differentiable with $L$‑Lipschitz gradient:
\[
\|\nabla f(x)-\nabla f(y)\|\le L\|x-y\|,\qquad\forall x,y\in\mathbb{R}^{d}.
\]
\end{assumption}
\begin{assumption}[Convexity of $g$]\label{ass:convex}
$g$ is proper, closed and convex (no smoothness required).
\end{assumption}
\begin{assumption}[Stepsize]\label{ass:stepsize}
The stepsize satisfies $0<\eta\le 1/L$.
\end{assumption}

\section*{Main Convergence Theorem}
\begin{theorem}[Ergodic $O(1/k)$ rate]\label{thm:conv}
Let $\{x_{k}\}$ be generated by the proximal gradient method under Assumptions \ref{ass:smooth}--\ref{ass:stepsize} and let $x^{\star}\in\arg\min F$. Then for any $k\ge1$,
\[
F(x_{k})-F(x^{\star})\le\frac{\|x_{0}-x^{\star}\|^{2}}{2\eta k}.
\]
Consequently,
\[
\min_{0\le i\le k-1}\bigl\{F(x_{i})-F(x^{\star})\bigr\}=O\!\left(\frac{1}{k}\right).
\]
\end{theorem}
\begin{theorem}[Linear rate under strong convexity]\label{thm:strong}
If, in addition, $F$ is $\mu$‑strongly convex ($\mu>0$) and $\eta\le 1/L$, then
\[
F(x_{k})-F(x^{\star})\le \bigl(1-\eta\mu\bigr)^{k}\,\bigl(F(x_{0})-F(x^{\star})\bigr).
\]
\end{theorem}

\section*{Theoretical Properties}
\begin{itemize}
\item \textbf{Stability:} The iteration is a contraction in the metric induced by the proximal operator because of its non‑expansiveness (Lemma~\ref{lem:nonexpansive}).
\item \textbf{Hyperparameter sensitivity:} The stepsize bound $\eta\le 1/L$ is tight; larger $\eta$ may violate the descent inequality.
\item \textbf{Comparison:} When $g\equiv0$, PGM reduces to gradient descent with the same $O(1/k)$ rate. When $g$ is an indicator of a convex set, PGM becomes projected gradient descent.
\end{itemize}

\section*{Preliminaries (Definitions and Lemmas)}
\begin{lemma}[Descent Lemma]\label{lem:descent}
If $\nabla f$ is $L$‑Lipschitz, then for any $x,y$,
\[
f(y)\le f(x)+\langle\nabla f(x),y-x\rangle+\frac{L}{2}\|y-x\|^{2}.
\]
\end{lemma}
\begin{lemma}[Non‑expansiveness of $\operatorname{prox}$]\label{lem:nonexpansive}
For any $u,v\in\mathbb{R}^{d}$ and any $\eta>0$,
\[
\bigl\|\operatorname{prox}_{\eta g}(u)-\operatorname{prox}_{\eta g}(v)\bigr\|
\le\|u-v\|.
\]
\end{lemma}
\begin{proof}
Let $p=\operatorname{prox}_{\eta g}(u)$ and $q=\operatorname{prox}_{\eta g}(v)$. By the optimality condition of the proximal map,
\[
\frac{1}{\eta}(u-p)\in\partial g(p),\qquad 
\frac{1}{\eta}(v-q)\in\partial g(q).
\]
Convexity of $g$ yields $\langle p-q,\,\frac{1}{\eta}(u-p)-\frac{1}{\eta}(v-q)\rangle\ge0$. Rearranging gives
\[
\|p-q\|^{2}\le\langle p-q,\,u-v\rangle\le\|p-q\|\|u-v\|,
\]
hence $\|p-q\|\le\|u-v\|$.
\end{proof}

\section*{Main Proof (Step‑by‑Step Derivation)}
We prove Theorem~\ref{thm:conv}.  
Let $x^{\star}\in\arg\min F$ and denote $v_{k}=x_{k}-\eta\nabla f(x_{k})$.

\begin{enumerate}[label=[Step \arabic*], leftmargin=*]
\item \textbf{Optimality condition of the proximal step.}  
By definition of the proximal operator,
\[
x_{k+1}= \operatorname{prox}_{\eta g}(v_{k})
\iff
\frac{1}{\eta}\bigl(v_{k}-x_{k+1}\bigr)\in\partial g(x_{k+1}).\tag{1}
\]

\item \textbf{Subgradient inequality for $g$.}  
For any $u$, convexity of $g$ gives
\[
g(u)\ge g(x_{k+1})+\Big\langle\frac{1}{\eta}(v_{k}-x_{k+1}),\,u-x_{k+1}\Big\rangle.\tag{2}
\]

\item \textbf{Apply (2) with $u=x^{\star}$.}  
\[
g(x^{\star})\ge g(x_{k+1})+\frac{1}{\eta}\langle v_{k}-x_{k+1},\,x^{\star}-x_{k+1}\rangle.\tag{3}
\]

\item \textbf{Expand the inner product.}  
Using $v_{k}=x_{k}-\eta\nabla f(x_{k})$,
\[
\langle v_{k}-x_{k+1},\,x^{\star}-x_{k+1}\rangle
= \langle x_{k}-x_{k+1},\,x^{\star}-x_{k+1}\rangle-\eta\langle\nabla f(x_{k}),\,x^{\star}-x_{k+1}\rangle.\tag{4}
\]

\item \textbf{Rewrite the first term of (4).}  
\[
\langle x_{k}-x_{k+1},\,x^{\star}-x_{k+1}\rangle
=\frac{1}{2}\Bigl(\|x_{k}-x^{\star}\|^{2}-\|x_{k+1}-x^{\star}\|^{2}-\|x_{k+1}-x_{k}\|^{2}\Bigr).\tag{5}
\]

\item \textbf{Insert (4)–(5) into (3).}  
Multiplying (3) by $\eta$ and using (4)–(5) yields
\[
\begin{aligned}
\eta\bigl(g(x^{\star})-g(x_{k+1})\bigr)
&\ge\frac{1}{2}\bigl(\|x_{k}-x^{\star}\|^{2}-\|x_{k+1}-x^{\star}\|^{2}-\|x_{k+1}-x_{k}\|^{2}\bigr)\\
&\quad-\eta\langle\nabla f(x_{k}),\,x^{\star}-x_{k+1}\rangle.\tag{6}
\end{aligned}
\]

\item \textbf{Add $f(x_{k+1})-f(x^{\star})$ to both sides of (6).}  
Using the descent lemma (Lemma~\ref{lem:descent}) with $y=x_{k+1}$ and $x=x_{k}$,
\[
f(x_{k+1})\le f(x_{k})+\langle\nabla f(x_{k}),\,x_{k+1}-x_{k}\rangle+\frac{L}{2}\|x_{k+1}-x_{k}\|^{2}.\tag{7}
\]
Rearrange (7) as
\[
-\langle\nabla f(x_{k}),\,x^{\star}-x_{k+1}\rangle
\le f(x_{k})-f(x^{\star})-\langle\nabla f(x_{k}),\,x_{k}-x^{\star}\rangle
+\frac{L}{2}\|x_{k+1}-x_{k}\|^{2}.\tag{8}
\]

\item \textbf{Combine (6) and (8).}  
Substituting (8) into (6) gives
\[
\begin{aligned}
F(x_{k+1})-F(x^{\star})
&\le \frac{1}{2\eta}\bigl(\|x_{k}-x^{\star}\|^{2}-\|x_{k+1}-x^{\star}\|^{2}\bigr)\\
&\quad+\Bigl(\frac{L}{2}-\frac{1}{2\eta}\Bigr)\|x_{k+1}-x_{k}\|^{2}.
\end{aligned}\tag{9}
\]

\item \textbf{Use the stepsize condition.}  
Assumption~\ref{ass:stepsize} guarantees $\eta\le 1/L$, i.e. $\frac{L}{2}-\frac{1}{2\eta}\le0$. Hence the second term in (9) is non‑positive and can be dropped:
\[
F(x_{k+1})-F(x^{\star})\le \frac{1}{2\eta}\bigl(\|x_{k}-x^{\star}\|^{2}-\|x_{k+1}-x^{\star}\|^{2}\bigr).\tag{10}
\]

\item \textbf{Telescoping sum.}  
Sum (10) for $k=0,\dots, K-1$:
\[
\sum_{k=0}^{K-1}\bigl(F(x_{k+1})-F(x^{\star})\bigr)
\le\frac{1}{2\eta}\bigl(\|x_{0}-x^{\star}\|^{2}-\|x_{K}-x^{\star}\|^{2}\bigr)
\le\frac{\|x_{0}-x^{\star}\|^{2}}{2\eta}.
\tag{11}
\]

\item \textbf{Derive the rate.}  
Since each term in the left‑hand side of (11) is non‑negative,
\[
K\bigl(F(x_{K})-F(x^{\star})\bigr)\le\frac{\|x_{0}-x^{\star}\|^{2}}{2\eta},
\]
which yields
\[
F(x_{K})-F(x^{\star})\le\frac{\|x_{0}-x^{\star}\|^{2}}{2\eta K}.
\tag{12}
\]
This is exactly the bound claimed in Theorem~\ref{thm:conv}.
\end{enumerate}

\section*{Rate Derivation (With Constants)}
From (12) we identify the explicit constant $C=\|x_{0}-x^{\star}\|^{2}/(2\eta)$. Hence
\[
F(x_{k})-F(x^{\star})\le\frac{C}{k},\qquad k\ge1.
\]

\section*{Special Cases}
\begin{enumerate}[label=\arabic*.]
\item \textbf{Pure gradient descent.} If $g\equiv0$, then $\operatorname{prox}_{\eta g}$ is the identity and the algorithm coincides with standard gradient descent; the same $O(1/k)$ bound holds.
\item \textbf{Projected gradient descent.} If $g$ is the indicator of a closed convex set $\mathcal{C}$, $\operatorname{prox}_{\eta g}$ becomes the Euclidean projection onto $\mathcal{C}$. The proof above reduces to the classic projected gradient analysis.
\item \textbf{Strongly convex $F$.} Adding $\mu$‑strong convexity ($F(y)\ge F(x)+\langle\nabla F(x),y-x\rangle+\frac{\mu}{2}\|y-x\|^{2}$) to the steps above yields a linear contraction:
\[
\|x_{k+1}-x^{\star}\|^{2}\le(1-\eta\mu)\|x_{k}-x^{\star}\|^{2},
\]
which directly gives the rate in Theorem~\ref{thm:strong}.
\end{enumerate}

\section*{Discussion}
The convergence proof hinges on two elementary properties:

\begin{itemize}
\item \textbf{Descent Lemma} for the smooth part $f$ (Lemma~\ref{lem:descent}).
\item \textbf{Non‑expansiveness} of the proximal map (Lemma~\ref{lem:nonexpansive}).
\end{itemize}
No stochasticity or variance assumptions are required; the analysis is deterministic and holds for any convex $g$. The stepsize bound $\eta\le1/L$ is both necessary and sufficient for the monotone decrease of $F$ and for the $O(1/k)$ ergodic convergence.

\end{document}